\section{Phase 3: BLOM-Strict --- Local Variational Definition and Well-Posedness}
\label{sec:phase3_blom_strict}

This section turns BLOM from an informal localization idea into a mathematically well-defined object.
The key task is to define a local variational problem on each window and to prove that the problem is well-posed, namely:
\begin{enumerate}
    \item the local admissible set is nonempty and affine;
    \item the local objective is strictly convex on that affine set;
    \item the induced local linear system is nonsingular;
    \item the local minimizer is unique.
\end{enumerate}
This theorem will play the role of the local analogue of the global MINCO mother theorem established in Phase~1.

Throughout this section, all notation is inherited from Phases~0--2.
In particular,
\[
\theta=(q_1,\dots,q_{M-1},T_1,\dots,T_M)^\top,
\qquad
T_i>0,
\qquad
\mathcal T=\sum_{i=1}^{M}T_i,
\]
and the canonical setting remains
\[
s=4,\qquad d_i=1
\]
unless otherwise stated.
For general statements we keep \(s\ge 2\) symbolic, and only specialize to \(s=4\) when needed.

\subsection{Local windows and local knot times}
\label{subsec:phase3_windows}

For each segment index \(i\in\{1,\dots,M\}\) and bandwidth parameter \(k\in\mathbb N\), define the window
\begin{equation}
\label{eq:phase3_window}
W(i,k)
=
\{\,i-\lfloor k/2\rfloor,\dots,i+\lfloor k/2\rfloor\,\}\cap[1,M].
\end{equation}
Let
\[
L_i:=\min W(i,k),\qquad R_i:=\max W(i,k),\qquad m_i:=|W(i,k)|=R_i-L_i+1.
\]
Thus the \(i\)-th local problem is posed on the segment set
\[
\{L_i,L_i+1,\dots,R_i\}.
\]

The corresponding local knot set is
\[
\Xi_{i,k}:=\{L_i-1,L_i,\dots,R_i\}.
\]
We introduce local cumulative times
\begin{equation}
\label{eq:phase3_local_knots}
\sigma_{L_i-1}:=0,
\qquad
\sigma_r:=\sum_{\ell=L_i}^{r}T_\ell,
\qquad
r=L_i,\dots,R_i,
\end{equation}
and denote the local horizon by
\begin{equation}
\label{eq:phase3_local_horizon}
\mathcal T_{i,k}:=\sigma_{R_i}=\sum_{\ell=L_i}^{R_i}T_\ell.
\end{equation}

\subsection{Local spline space and affine constraint set}
\label{subsec:phase3_local_space_constraints}

The local BLOM-Strict problem is defined on a finite-dimensional spline space.
Let
\[
\mathbb P_{2s-1}
\]
denote the space of real polynomials of degree at most \(2s-1\).
We define the local \(C^{s-1}\) spline space
\begin{equation}
\label{eq:phase3_local_spline_space}
\mathcal S_{i,k}^{(s-1)}(T)
:=
\left\{
u=(u_{L_i},\dots,u_{R_i})
\ \middle|\
\begin{aligned}
& u_j\in \mathbb P_{2s-1},\qquad j=L_i,\dots,R_i,\\
& u_j^{(\ell)}(T_j)=u_{j+1}^{(\ell)}(0),\\
& \hspace{1.8cm}\ell=0,\dots,s-1,\quad j=L_i,\dots,R_i-1
\end{aligned}
\right\}.
\end{equation}
Thus the segmentwise functions are polynomial of degree at most \(2s-1\) and are \(C^{s-1}\) across all internal knots of the window.

\paragraph{Interpolation constraints.}
For every \(u\in\mathcal S_{i,k}^{(s-1)}(T)\), we impose exact waypoint interpolation at all window knots:
\begin{equation}
\label{eq:phase3_window_interpolation}
u(\sigma_r)=q_r,
\qquad r\in \Xi_{i,k}=\{L_i-1,\dots,R_i\}.
\end{equation}

\paragraph{Physical boundary inheritance.}
If the local window touches a true global boundary, then BLOM-Strict inherits the physical boundary jet from Phase~1.
Specifically:
\begin{itemize}
    \item if \(L_i=1\), then in addition to \(u(0)=q_0\) already enforced by \eqref{eq:phase3_window_interpolation}, we require
    \begin{equation}
    \label{eq:phase3_left_physical_jet}
    u^{(\ell)}(0)=\zeta^-_\ell,
    \qquad \ell=1,\dots,s-1;
    \end{equation}
    \item if \(R_i=M\), then in addition to \(u(\mathcal T_{i,k})=q_M\), we require
    \begin{equation}
    \label{eq:phase3_right_physical_jet}
    u^{(\ell)}(\mathcal T_{i,k})=\zeta^+_\ell,
    \qquad \ell=1,\dots,s-1.
    \end{equation}
\end{itemize}

No derivative constraints are imposed at artificial window boundaries.
The corresponding natural boundary conditions will instead be derived from first-order optimality.

\begin{definition}[Local admissible set]
\label{def:phase3_local_admissible}
The BLOM-Strict admissible set for the window \(W(i,k)\) is
\begin{equation}
\label{eq:phase3_local_admissible}
\mathcal A_{i,k}(q,T;\zeta^-,\zeta^+)
:=
\left\{
u\in\mathcal S_{i,k}^{(s-1)}(T)
\ \middle|\
\eqref{eq:phase3_window_interpolation}\text{ holds, and }
\eqref{eq:phase3_left_physical_jet},\eqref{eq:phase3_right_physical_jet}
\text{ are imposed whenever applicable}
\right\}.
\end{equation}
\end{definition}

\begin{remark}[Why the window space is only \(C^{s-1}\) a priori]
\label{rem:phase3_why_c_s_minus_1}
The feasible space in \eqref{eq:phase3_local_spline_space} imposes only the continuity that is structurally required by the \(s\)-fold integrator chain, namely \(C^{s-1}\).
The higher continuity \(C^{2s-2}\) will be derived from the optimality conditions, exactly as in the global MINCO mother theorem of Phase~1.
This is essential: if one imposed \(C^{2s-2}\) and natural boundary conditions as hard constraints from the start, then the local variational problem would lose the precise analogy with the global MINCO construction.
\end{remark}

\subsection{The BLOM-Strict local variational problem}
\label{subsec:phase3_local_problem}

For each window \(W(i,k)\), define the local control effort functional
\begin{equation}
\label{eq:phase3_local_energy}
J_{i,k}[u]
:=
\sum_{j=L_i}^{R_i}\int_{0}^{T_j}\bigl|u_j^{(s)}(\tau)\bigr|^2\,d\tau.
\end{equation}

\begin{definition}[BLOM-Strict local problem]
\label{def:phase3_blom_strict_problem}
The BLOM-Strict local problem on the window \(W(i,k)\) is
\begin{equation}
\label{eq:phase3_local_problem}
\boxed{
\min_{u\in\mathcal A_{i,k}(q,T;\zeta^-,\zeta^+)}
J_{i,k}[u].
}
\end{equation}
Its unique minimizer, when it exists, is denoted by
\[
u_{i,k}^\star.
\]
\end{definition}

\subsection{Natural boundary conditions at artificial window boundaries}
\label{subsec:phase3_natural_bc}

We now derive the exact natural boundary conditions induced by \eqref{eq:phase3_local_problem}.

\begin{proposition}[Natural boundary conditions for BLOM-Strict]
\label{prop:phase3_natural_bc}
Let \(u_{i,k}^\star\) be a local minimizer of \eqref{eq:phase3_local_problem}.
Then:
\begin{itemize}
    \item if \(L_i>1\), i.e.\ the left boundary of the window is artificial, then
    \begin{equation}
    \label{eq:phase3_left_natural_bc}
    \bigl(u_{i,k}^\star\bigr)^{(r)}(0)=0,
    \qquad r=s,\dots,2s-2;
    \end{equation}
    \item if \(R_i<M\), i.e.\ the right boundary of the window is artificial, then
    \begin{equation}
    \label{eq:phase3_right_natural_bc}
    \bigl(u_{i,k}^\star\bigr)^{(r)}(\mathcal T_{i,k})=0,
    \qquad r=s,\dots,2s-2.
    \end{equation}
\end{itemize}
\end{proposition}

\begin{proof}
We prove the left boundary statement; the right one is identical.

Assume \(L_i>1\), so the left boundary is artificial.
Let \(h\in \mathcal S_{i,k}^{(s-1)}(T)\) be any feasible variation of \(u_{i,k}^\star\), i.e.\ \(h\) satisfies the homogeneous form of all affine constraints.
In particular,
\[
h(0)=0
\]
because the left endpoint value is fixed by interpolation, while
\[
h^{(\ell)}(0),\qquad \ell=1,\dots,s-1,
\]
are free because no derivative constraints are imposed at the artificial boundary.

Since \(u_{i,k}^\star\) is a minimizer, the first variation vanishes:
\[
0
=
\delta J_{i,k}[u_{i,k}^\star;h]
=
2\sum_{j=L_i}^{R_i}\int_0^{T_j}
(u_{i,k}^\star)_j^{(s)}(\tau)\, h_j^{(s)}(\tau)\,d\tau.
\]
Integrating by parts \(s\) times on the first segment and collecting the left boundary terms gives
\[
\sum_{r=0}^{s-1}
(-1)^r
\bigl(u_{i,k}^\star\bigr)^{(s+r)}(0)\,
h^{(s-1-r)}(0).
\]
Because \(h(0)=0\), the term corresponding to \(r=s-1\) vanishes automatically.
The remaining traces
\[
h^{(1)}(0),\dots,h^{(s-1)}(0)
\]
are free, hence the coefficients in front of them must vanish:
\[
\bigl(u_{i,k}^\star\bigr)^{(s)}(0)=\cdots=
\bigl(u_{i,k}^\star\bigr)^{(2s-2)}(0)=0.
\]
This is exactly \eqref{eq:phase3_left_natural_bc}.
\end{proof}

\subsection{Window admissibility and the local nullspace obstruction}
\label{subsec:phase3_window_admissibility}

The strict convexity of the local problem depends on whether the zero-cost polynomial nullspace is killed by the local constraints.

\begin{definition}[(\(s,k\))-admissible window]
\label{def:phase3_window_admissible}
A window \(W(i,k)\) is called \emph{(\(s,k\))-admissible} if either
\begin{enumerate}
    \item it touches at least one physical global boundary, i.e.\ \(L_i=1\) or \(R_i=M\); or
    \item it is a purely interior window and
    \begin{equation}
    \label{eq:phase3_admissibility_condition}
    |W(i,k)|=m_i\ge s-1.
    \end{equation}
\end{enumerate}
\end{definition}

\begin{remark}[Canonical specialization]
\label{rem:phase3_canonical_admissibility}
In the canonical setting \(s=4,\ k=2\), every interior window has size
\[
|W(i,2)|=3=s-1,
\]
so condition \eqref{eq:phase3_admissibility_condition} holds exactly at the threshold.
Boundary windows are admissible by physical boundary inheritance.
\end{remark}

The reason for Definition~\ref{def:phase3_window_admissible} is the following nullspace lemma.

\begin{lemma}[Triviality of the zero-cost nullspace]
\label{lem:phase3_nullspace}
Assume that \(W(i,k)\) is \((s,k)\)-admissible.
Let \(h\in \mathcal A_{i,k}(0,T;0,0)\) be a homogeneous feasible direction, i.e.\ all interpolation data and all inherited boundary jet data are zero.
If
\begin{equation}
\label{eq:phase3_zero_cost}
J_{i,k}[h]=0,
\end{equation}
then \(h\equiv 0\) on the whole window.
\end{lemma}

\begin{proof}
Condition \eqref{eq:phase3_zero_cost} implies
\[
h_j^{(s)}(\tau)=0
\qquad
\text{for every } j=L_i,\dots,R_i \text{ and all } \tau\in[0,T_j].
\]
Hence every segment \(h_j\) is a polynomial of degree at most \(s-1\).

Since \(h\in \mathcal S_{i,k}^{(s-1)}(T)\), the derivatives of orders \(0,\dots,s-1\) match at every internal knot of the window.
Therefore all segment polynomials are restrictions of one and the same global polynomial
\[
\pi\in \mathbb P_{s-1}
\]
defined on the local interval \([0,\mathcal T_{i,k}]\).

We distinguish two cases.

\paragraph{Case 1: the window touches a physical boundary.}
Suppose \(L_i=1\).
Then the homogeneous version of \eqref{eq:phase3_window_interpolation} and \eqref{eq:phase3_left_physical_jet} gives
\[
\pi^{(\ell)}(0)=0,
\qquad
\ell=0,1,\dots,s-1.
\]
Hence all coefficients of the degree-\((s-1)\) polynomial \(\pi\) vanish, so \(\pi\equiv 0\), and thus \(h\equiv 0\).
The case \(R_i=M\) is identical.

\paragraph{Case 2: the window is interior.}
Then \(L_i>1\) and \(R_i<M\), and by admissibility \(m_i\ge s-1\).
The homogeneous interpolation constraints imply that
\[
\pi(\sigma_r)=0,
\qquad
r=L_i-1,\dots,R_i.
\]
Thus \(\pi\) has \(m_i+1\) distinct roots.
Since \(m_i+1\ge s\) and \(\deg \pi\le s-1\), it follows that \(\pi\equiv 0\).
Therefore \(h\equiv 0\).
\end{proof}

\subsection{Theorem 2: existence, uniqueness, and optimality structure of BLOM-Strict}
\label{subsec:phase3_theorem2}

We can now state the main theorem of this phase.

\begin{theorem}[BLOM-Strict existence, uniqueness, and local optimality structure]
\label{thm:phase3_blom_strict_wellposedness}
Assume that \(T_j>0\) for all \(j\in W(i,k)\) and that the window \(W(i,k)\) is \((s,k)\)-admissible in the sense of Definition~\ref{def:phase3_window_admissible}.
Then the BLOM-Strict local problem \eqref{eq:phase3_local_problem} admits a unique minimizer
\[
u_{i,k}^\star \in \mathcal A_{i,k}(q,T;\zeta^-,\zeta^+).
\]
Moreover, the minimizer satisfies the following properties:

\begin{enumerate}
    \item \textbf{Polynomial structure.}
    Each segment of \(u_{i,k}^\star\) is a polynomial of degree at most \(2s-1\) by construction.

    \item \textbf{Interpolation and inherited boundary data.}
    The minimizer satisfies all interpolation constraints \eqref{eq:phase3_window_interpolation}, and whenever the window touches a physical boundary it satisfies the inherited boundary jet conditions \eqref{eq:phase3_left_physical_jet} and/or \eqref{eq:phase3_right_physical_jet}.

    \item \textbf{Natural boundary conditions at artificial boundaries.}
    At every artificial window boundary, the natural conditions
    \eqref{eq:phase3_left_natural_bc}--\eqref{eq:phase3_right_natural_bc} hold.

    \item \textbf{Higher continuity inside the window.}
    The minimizer is \(C^{2s-2}\) across every internal knot of the window, i.e.
    \begin{equation}
    \label{eq:phase3_high_continuity}
    (u_{i,k}^\star)^{(r)}(\sigma_\ell^-)
    =
    (u_{i,k}^\star)^{(r)}(\sigma_\ell^+),
    \qquad
    r=0,\dots,2s-2,
    \qquad
    \ell=L_i,\dots,R_i-1.
    \end{equation}
\end{enumerate}

Equivalently, in coefficient coordinates the problem is a strictly convex equality-constrained quadratic program whose reduced local linear system is nonsingular.
\end{theorem}

\begin{proof}
The proof is divided into four steps.

\paragraph{Step 1: nonemptiness of the admissible set.}
We explicitly construct a feasible element.
Choose derivative data of orders \(1,\dots,s-1\) at every internal knot inside the window to be zero.
If the window touches a physical boundary, keep the prescribed boundary jet there; if not, set the artificial boundary derivatives of orders \(1,\dots,s-1\) to zero.
Then, for each segment \(j\in W(i,k)\), the two endpoint values and the two endpoint derivative jets of orders \(1,\dots,s-1\) determine a unique Hermite polynomial of degree at most \(2s-1\).
By construction, the resulting piecewise polynomial belongs to \(\mathcal S_{i,k}^{(s-1)}(T)\) and satisfies all affine constraints.
Hence
\[
\mathcal A_{i,k}(q,T;\zeta^-,\zeta^+)\neq\varnothing.
\]

\paragraph{Step 2: strict convexity on the affine feasible set.}
Let \(u,v\in \mathcal A_{i,k}(q,T;\zeta^-,\zeta^+)\), and let
\[
h:=u-v\in \mathcal A_{i,k}(0,T;0,0)
\]
be the corresponding homogeneous feasible direction.
Since \(J_{i,k}\) is quadratic,
\[
J_{i,k}\!\left(\frac{u+v}{2}\right)
=
\frac{1}{2}J_{i,k}[u]+\frac{1}{2}J_{i,k}[v]-\frac{1}{4}J_{i,k}[h].
\]
Therefore strict convexity on the affine feasible set will follow once we prove:
\[
h\neq 0 \quad\Longrightarrow\quad J_{i,k}[h]>0.
\]
But if \(J_{i,k}[h]=0\), then Lemma~\ref{lem:phase3_nullspace} yields \(h\equiv 0\), contradiction.
Hence \(J_{i,k}\) is strictly convex on \(\mathcal A_{i,k}(q,T;\zeta^-,\zeta^+)\).

\paragraph{Step 3: existence and uniqueness.}
Because \(\mathcal S_{i,k}^{(s-1)}(T)\) is finite-dimensional, the admissible set
\[
\mathcal A_{i,k}(q,T;\zeta^-,\zeta^+)
\]
is a finite-dimensional affine subspace.
The functional \(J_{i,k}\) is continuous and strictly convex on this affine subspace.
Hence it admits a unique minimizer \(u_{i,k}^\star\).

\paragraph{Step 4: optimality conditions.}
Let \(h\in \mathcal A_{i,k}(0,T;0,0)\) be any homogeneous feasible variation.
Then the first variation vanishes:
\[
0
=
\delta J_{i,k}[u_{i,k}^\star;h]
=
2\sum_{j=L_i}^{R_i}\int_{0}^{T_j}
(u_{i,k}^\star)_j^{(s)}(\tau)\,h_j^{(s)}(\tau)\,d\tau.
\]
Integrating by parts \(s\) times on each segment gives boundary terms at the internal knots and, if present, at artificial window boundaries.

At an artificial window boundary, Proposition~\ref{prop:phase3_natural_bc} follows because
\(h\) has fixed value but free derivatives of orders \(1,\dots,s-1\) there.

At an internal knot \(\sigma_\ell\), the traces
\[
h^{(1)}(\sigma_\ell),\dots,h^{(s-1)}(\sigma_\ell)
\]
are arbitrary, because the feasible space imposes only \(C^{s-1}\) continuity.
The knot terms obtained by integration by parts therefore force the jumps of the higher derivatives to vanish:
\[
\llbracket (u_{i,k}^\star)^{(s)} \rrbracket(\sigma_\ell)
=
\cdots
=
\llbracket (u_{i,k}^\star)^{(2s-2)} \rrbracket(\sigma_\ell)
=0.
\]
Since \(u_{i,k}^\star\in \mathcal S_{i,k}^{(s-1)}(T)\) already has continuity of orders \(0,\dots,s-1\), this proves \eqref{eq:phase3_high_continuity}.
The remaining claims are immediate from the definition of the admissible set.
\end{proof}

\subsection{Coefficient-space formulation and nonsingularity of the reduced local system}
\label{subsec:phase3_coefficient_qp}

Theorem~\ref{thm:phase3_blom_strict_wellposedness} implies that the local problem is a well-posed finite-dimensional quadratic program.

For each segment \(j\in W(i,k)\), write
\[
u_j(\tau)=c_j^\top \beta_s(\tau),
\qquad
c_j\in\mathbb R^{2s},
\qquad
\beta_s(\tau)=(1,\tau,\dots,\tau^{2s-1})^\top.
\]
Stack the local coefficients into
\[
c_{i,k}^{\mathrm{loc}}
=
(c_{L_i}^\top,\dots,c_{R_i}^\top)^\top
\in\mathbb R^{2s m_i}.
\]
Then the local objective takes the quadratic form
\begin{equation}
\label{eq:phase3_local_quadratic_form}
J_{i,k}[u]
=
\frac{1}{2}\,(c_{i,k}^{\mathrm{loc}})^\top
H_{i,k}(T)\,
c_{i,k}^{\mathrm{loc}},
\end{equation}
where
\begin{equation}
\label{eq:phase3_local_hessian}
H_{i,k}(T)
=
\operatorname{blkdiag}\!\bigl(Q_s(T_{L_i}),\dots,Q_s(T_{R_i})\bigr),
\end{equation}
and the per-segment Gram block is
\begin{equation}
\label{eq:phase3_gram_block}
(Q_s(T))_{\alpha\beta}
=
\int_0^T
\frac{d^s}{d\tau^s}\tau^{\alpha-1}\,
\frac{d^s}{d\tau^s}\tau^{\beta-1}\,d\tau,
\qquad
\alpha,\beta=1,\dots,2s.
\end{equation}

All affine constraints can be written compactly as
\begin{equation}
\label{eq:phase3_local_constraints_matrix}
G_{i,k}(T)\,c_{i,k}^{\mathrm{loc}}
=
d_{i,k}(q,T;\zeta^-,\zeta^+).
\end{equation}
Hence the local BLOM-Strict problem is
\begin{equation}
\label{eq:phase3_local_qp}
\boxed{
\min_{c_{i,k}^{\mathrm{loc}}\in\mathbb R^{2s m_i}}
\frac{1}{2}\,(c_{i,k}^{\mathrm{loc}})^\top H_{i,k}(T)c_{i,k}^{\mathrm{loc}}
\quad\text{s.t.}\quad
G_{i,k}(T)c_{i,k}^{\mathrm{loc}}
=
d_{i,k}(q,T;\zeta^-,\zeta^+).
}
\end{equation}

\begin{proposition}[Nonsingularity of the reduced local system]
\label{prop:phase3_reduced_system}
Assume the hypotheses of Theorem~\ref{thm:phase3_blom_strict_wellposedness}.
Let \(c_p\) be any particular feasible coefficient vector satisfying
\[
G_{i,k}(T)c_p=d_{i,k}(q,T;\zeta^-,\zeta^+),
\]
and let \(N_{i,k}(T)\) be a basis matrix whose columns span \(\ker G_{i,k}(T)\).
Then every feasible coefficient vector can be written as
\[
c_{i,k}^{\mathrm{loc}}=c_p+N_{i,k}(T)z,
\]
and the reduced Hessian
\begin{equation}
\label{eq:phase3_reduced_hessian}
R_{i,k}(T):=
N_{i,k}(T)^\top H_{i,k}(T)N_{i,k}(T)
\end{equation}
is symmetric positive definite.
Consequently, the reduced local linear system
\begin{equation}
\label{eq:phase3_reduced_linear_system}
R_{i,k}(T)z^\star
=
-
N_{i,k}(T)^\top H_{i,k}(T)c_p
\end{equation}
is nonsingular and determines the unique local minimizer.
\end{proposition}

\begin{proof}
Symmetry of \(R_{i,k}(T)\) is immediate from that of \(H_{i,k}(T)\).

Let \(z\neq 0\), and set
\[
\delta c:=N_{i,k}(T)z\in \ker G_{i,k}(T).
\]
Let \(\delta u\) be the corresponding homogeneous feasible local spline.
Then
\[
z^\top R_{i,k}(T)z
=
(\delta c)^\top H_{i,k}(T)\delta c
=
2J_{i,k}[\delta u].
\]
If \(z^\top R_{i,k}(T)z=0\), then \(J_{i,k}[\delta u]=0\), and Lemma~\ref{lem:phase3_nullspace} implies \(\delta u\equiv 0\).
Hence \(\delta c=0\), and since the columns of \(N_{i,k}(T)\) are linearly independent, it follows that \(z=0\), contradiction.
Therefore \(R_{i,k}(T)\) is positive definite and in particular nonsingular.
The reduced system \eqref{eq:phase3_reduced_linear_system} therefore has a unique solution \(z^\star\), and hence a unique feasible minimizer \(c_p+N_{i,k}(T)z^\star\).
\end{proof}

\subsection{Definition of the BLOM trajectory class}
\label{subsec:phase3_blom_trajectory_class}

We can now define BLOM as a parameter-to-segment map, and therefore as a trajectory class.

For each segment index \(i\), let
\[
u_{i,k}^\star
\]
be the unique local minimizer from Theorem~\ref{thm:phase3_blom_strict_wellposedness}, and let
\[
c_{i,k}^{\mathrm{loc},\star}
=
(c_{L_i}^\top,\dots,c_i^\top,\dots,c_{R_i}^\top)^\top
\]
be its local coefficient vector.
Define the segment extraction operator
\[
\Pi_i:\mathbb R^{2s m_i}\to\mathbb R^{2s}
\]
that extracts the block corresponding to the central segment \(i\).
The \(i\)-th BLOM-Strict segment coefficient is then
\begin{equation}
\label{eq:phase3_blom_segment_coeff}
c_i^{\mathrm{BLOM},k}(q,T;\zeta^-,\zeta^+)
:=
\Pi_i\, c_{i,k}^{\mathrm{loc},\star}(q,T;\zeta^-,\zeta^+).
\end{equation}
The corresponding segment polynomial is
\begin{equation}
\label{eq:phase3_blom_segment_poly}
p_i^{\mathrm{BLOM},k}(\tau)
=
\bigl(c_i^{\mathrm{BLOM},k}(q,T;\zeta^-,\zeta^+)\bigr)^\top \beta_s(\tau),
\qquad
\tau\in[0,T_i].
\end{equation}

\begin{definition}[BLOM trajectory class]
\label{def:phase3_blom_class}
For fixed \(k\), define
\begin{equation}
\label{eq:phase3_blom_class}
\mathcal T_{\mathrm{BLOM},k}
:=
\left\{
p=(p_1,\dots,p_M)
\ \middle|\
\exists\,(q,T,\zeta^-,\zeta^+)
\text{ such that }
p_i=p_i^{\mathrm{BLOM},k}(q,T;\zeta^-,\zeta^+)
\ \forall i=1,\dots,M
\right\}.
\end{equation}
\end{definition}

\begin{remark}[What \(\mathcal T_{\mathrm{BLOM},k}\) means at this stage]
\label{rem:phase3_blom_class_scope}
Definition~\ref{def:phase3_blom_class} makes BLOM a mathematically precise map
\[
(q,T,\zeta^-,\zeta^+)\longmapsto (c_1^{\mathrm{BLOM},k},\dots,c_M^{\mathrm{BLOM},k}),
\]
and hence a legitimate trajectory class.
However, at this stage we do \emph{not} claim that neighboring segments extracted from different local windows automatically enjoy global \(C^{2s-2}\) compatibility across the full trajectory.
That issue is postponed to the later phase on global assembly and boundary consistency.
Phase~3 only guarantees that each local window problem is well-defined and uniquely solvable.
\end{remark}

\subsection{Canonical specialization to the minimum-snap three-point case}
\label{subsec:phase3_canonical_specialization}

In the canonical setting of this work,
\[
s=4,\qquad k=2.
\]
Hence each interior window contains three segments:
\[
W(i,2)=\{i-1,i,i+1\},
\qquad 2\le i\le M-1,
\]
and therefore
\[
|W(i,2)|=3=s-1.
\]
The admissibility condition of Definition~\ref{def:phase3_window_admissible} is satisfied at equality.
Accordingly:
\begin{itemize}
    \item every interior BLOM-Strict local problem has a unique minimizer because the zero-cost nullspace consists of global cubics, and four window knots eliminate that nullspace;
    \item boundary windows are uniquely solvable because they inherit the physical boundary jet from Phase~1.
\end{itemize}
Thus the canonical \(s=4,\ k=2\) BLOM-Strict construction is a mathematically well-posed object.

\subsection{Logical role of Phase~3 in the full BLOM theory}
\label{subsec:phase3_logical_role}

Phase~2 proved that exact global optimality and exact finite local support are incompatible under the sparse waypoint parameterization of Phase~0/1.
Phase~3 now provides the first rigorous replacement:
instead of the exact global coefficient map
\[
(q,T)\longmapsto c^\star(q,T),
\]
we define a family of uniquely solvable local variational problems
\[
(q,T)\longmapsto c_i^{\mathrm{BLOM},k}(q,T).
\]
Therefore BLOM-Strict is not a heuristic interpolation rule, but a well-posed local optimal mapping.
The next phase will convert this abstract local variational definition into an explicit analytic local system of the form
\[
A_k(T)x=B_k(q,T),
\]
which is the computational core needed for efficient evaluation and differentiation.